% ======================================================================================
% Nom     : reports/latex/sections/02_contexte_problematique.tex
% Rôle    : Section présentant le contexte et la problématique du projet.
% Auteur  : Maxime BRONNY
% Version : V1
% Cadre   : Réalisé dans le cadre du cours Techniques d'apprentissage artificiel
%           Master 1 Informatique Big Data
% Usage   : Fichier inclus dans main.tex lors de la compilation du rapport.
% ======================================================================================

\section{Contexte et problématique}
\label{sec:contexte}

\subsection{Le problème métier}

Lorsqu'une banque accorde un crédit, deux types d'erreurs sont possibles, et
leurs coûts sont très asymétriques. Si la banque prête à un client qui ne
remboursera pas (défaut non détecté), elle perd tout ou partie du capital
prêté ; si elle refuse un client qui aurait remboursé, elle ne perd qu'un manque
à gagner, les intérêts du prêt. Un bon système de prédiction du risque doit donc
non seulement être globalement précis, mais surtout détecter le plus possible de
futurs défauts, sans pour autant rejeter massivement les bons dossiers.

Le volume des demandes de crédit rend l'examen entièrement manuel impraticable,
tandis que les banques disposent d'historiques de remboursement riches et
structurés. Ce contexte, caractérisé par beaucoup de données étiquetées et une
décision binaire à automatiser, constitue un cadre idéal pour l'apprentissage
supervisé.

\subsection{Formulation en problème d'apprentissage}

Le problème se formule comme une \textbf{classification binaire supervisée} :
à partir d'un vecteur de caractéristiques décrivant l'emprunteur et son prêt
(âge, revenu, montant, taux d'intérêt, historique de crédit, etc.), prédire la
variable cible \texttt{loan\_status} :
\begin{itemize}
  \item classe $0$ : remboursement normal (environ $78\,\%$ des observations) ;
  \item classe $1$ : défaut de paiement (environ $22\,\%$ des observations).
\end{itemize}

Ce déséquilibre des classes est la difficulté centrale du problème : un modèle
naïf qui prédirait systématiquement « pas de défaut » atteindrait déjà environ
$78\,\%$ d'exactitude tout en étant parfaitement inutile, puisqu'il ne
détecterait aucun défaut. L'évaluation devra donc reposer sur des métriques
adaptées (précision, rappel, F1-score, courbe ROC), présentées en
section~\ref{sec:methodologie}.

\subsection{Problématique}

La problématique retenue pour ce projet est la suivante :

\begin{quote}
\itshape
À partir des caractéristiques d'un emprunteur et de son prêt, peut-on prédire
automatiquement le risque de défaut de paiement, et quelle méthode
d'apprentissage supervisé (une méthode linéaire, une méthode à base de règles ou
une méthode à base d'instances) offre le meilleur compromis entre la détection
des défauts et les erreurs de classification ?
\end{quote}

Répondre à cette question suppose un protocole expérimental rigoureux : mêmes
données, même prétraitement et même méthode de sélection des hyperparamètres
pour les trois modèles, et une évaluation finale sur des données jamais
utilisées pendant la conception.
